% !TeX program = xelatex
% UTF-8. Standalone document; compile from this directory or the repository root.
% Existing six-page transcription retained from the earlier Google Drive source set.
% pic01.jpg = 20261016_212038277.jpg (1x6B7dsRehpmYGTeYL1aLD8J2aYP8lcV0)
% pic02.jpg = 20261016_212046217.jpg (1glYNOk7ottTw6Ej2IAZVMZRwrOj-eB87)
% Original transcription page 3: 20261016_212053063.jpg (1z_0UpSsRHw0vrpzKEZlr1EwW8E5g4LU0)
% Original transcription page 4: 20261016_212101490.jpg (1iVh1lV3yWl-6HVbHM5Aw7GyFgEiTWAa2)
% pic05.jpg = 20261016_212108170.jpg (1V2w1zbVXoqqb0iN1VErXsSoAwigCY9iY)
% pic06.jpg = 20261016_212113932.jpg (14GBDSH9TcDvFrvt4FOp2lyIsdVToHrpu)
% Added two-page summary: Google Drive / Notebook / JIA.
% pic03.jpg = 20261018_192809528.jpg (1pjQgdJDESkfhVugMo1veFLs6VpYCkVWm)
% pic04.jpg = 20261018_192814024.jpg (1NftOpJpTbjZg4j3UplDPIRTw70ZxlBrD)
% These replace the previously occupied pic03/pic04 filenames as requested.
\documentclass[11pt,a4paper]{article}
\usepackage[margin=20mm]{geometry}
\usepackage{amsmath,amssymb}
\usepackage{graphicx}
\usepackage{longtable,array}
\usepackage{enumitem}
\usepackage[normalem]{ulem}
\usepackage{xepersian}
\settextfont{Amiri}
\setlatintextfont{DejaVu Serif}
\setlength{\parindent}{0pt}
\setlength{\parskip}{2pt}
\setlength{\emergencystretch}{3em}
\setlist[itemize]{itemsep=1pt,topsep=2pt}
\newcommand{\en}[1]{\lr{#1}}
\newcommand{\rng}[2]{\lr{#1--#2}}
\newcommand{\up}{\ensuremath{\uparrow}}
\newcommand{\dn}{\ensuremath{\downarrow}}
\newcommand{\unclear}[1]{[خوانش نامطمئن: #1]}
\newcommand{\pagehead}[1]{\section*{صفحهٔ دست‌نویس #1}}
\begin{document}
\begin{center}
{\Large\bfseries \en{146 -- JIA}}\\[3pt]
\en{Nelson [22]}
\end{center}

\pagehead{۱}
\begin{itemize}
\item شایع‌ترین بیماری روماتولوژیک در کودکان.
\item اختلالات خودایمنی که درگیری اصلی‌شان آرتریت است:
\item \en{JIA}: اتیولوژی و پاتوژنز آن به‌طور عمده ناشناخته است.
\item قبلاً به آن آرتریت روماتوئید جوانان می‌گفتند؛ الان آرتریت ایدیوپاتیک جوانان.
\end{itemize}

\subsection*{اپیدمیولوژی}
\begin{itemize}
\item در هر \en{100,000} کودک، \rng{0.8}{22.6} در سال رخ می‌دهد.
\item شیوع آن: \rng{7}{401} در هر \en{100,000} کودک.
\end{itemize}
\textbf{\en{Subtypes}:}
\begin{center}
\begin{tabular}{|l|c|c|}
\hline
نوع & فراوانی & نسبت \en{F/M}\\\hline
\en{Oligoarthritis} & شایع‌ترین؛ \rng{40}{50}\% & \en{3/1}\\\hline
\en{Polyarthritis} & \rng{25}{30}\% & \en{5/1}\\\hline
\en{Systemic JIA} & \rng{5}{15}\% & \en{1/1}\\\hline
\end{tabular}
\end{center}
\textbf{سن شایع بروز بیماری:}
\begin{itemize}
\item \en{Olygo}: \en{2--4 yo}.
\item \en{Poly}: \rng{2}{4} و \rng{10}{14} سال.
\item \en{Sys}: \rng{1}{5} سال.
\end{itemize}

\subsection*{اتیولوژی}
\begin{itemize}
\item استعداد ایمونوژنتیک؛ زمینهٔ خانوادگی.
\item برخی واریانت‌های \en{MHC} کلاس \en{I} و \en{II}.
\item ژن‌ها مرتبط با \en{IL-6}، \en{IL1-$\alpha$}، \en{TNF-$\alpha$}، \en{PTPN22}.
\end{itemize}

\subsection*{پاتوژنز}
\begin{itemize}
\item یک بیماری خودایمنی مرتبط با تغییر در ایمنی سلولی و هومورال.
\item لنفوسیت‌های \en{T} نقش مهمی دارند.
\item در \en{sJIA}: اختلال تنظیم ایمنی ذاتی + کمبود سلول‌های \en{T} اتوریاکتیو و اتوآنتی‌بادی‌ها.
\item به همین دلیل شاید بهتر باشد بگوییم \en{sJIA} یک اختلال خودالتهابی است که زمانی که یک آرتریت ایجاد می‌شود تازه خودایمنی می‌شود.
\item نقش \en{IL-1} بسیار مهم است.
\item سطح \en{IL-18} در \en{JIA} به شدت بالا می‌رود.
\end{itemize}

\subsection*{تظاهرات بالینی}
\begin{itemize}
\item وجود آرتریت \en{$\geq 6$} هفته برای تشخیص ضروری است.
\item آرتریت: التهاب و تورم داخل مفصلی یا وجود حداقل ۲ صفت: محدودیت \en{ROM}، تندرنس یا درد حین حرکت، گرمی مفصل.
\end{itemize}

\clearpage
\pagehead{۲}
\begin{itemize}
\item علائم اولیه معمولاً خفیف و شامل خشکی صبحگاهی، خستگی زودرس، کیفیت خواب کم.
\item مفاصل معمولاً اریتماتو نیستند.
\item آرتریت مفاصل بزرگ خصوصاً زانو در ابتدا رشد خطی را افزایش می‌دهد و منجر به بلندتر بودن اندام درگیر نسبت به اندام مقابل می‌شود.
\item التهاب با واسطهٔ واسکولر، بسته شدن زودهنگام صفحهٔ رشد می‌شود و در نهایت اندام درگیر نسبت به اندام مقابل کوتاه‌تر خواهد بود.
\end{itemize}

\subsection*{اولیگوآرتریت}
\begin{itemize}
\item حداکثر ۴ مفصل در ۶ ماه نخست بیماری درگیر می‌شوند.
\item اغلب یک مفصل درگیر می‌شود.
\item به‌طور عمده مفاصل بزرگ اندام تحتانی درگیر می‌شوند: زانو، مچ پا.
\item پس از ۶ ماه اگر تعداد مفاصل درگیر بیشتر از ۴ نشود: \en{Persistent Olygo}.
\item بشود: \en{Extended Oligoarthritis JIA}.
\item \en{Extended} پروگنوز بدتری دارد.
\end{itemize}

\subsection*{ویژگی‌های \en{JIA}}
\begin{itemize}
\item درگیری ایزولهٔ هیپ در \en{JIA} هرگز تظاهر اولیه نیست، بلکه مطرح‌کنندهٔ \en{ERA} است.
\item مثبت شدن \en{ANA} خطر یووئیت قدامی را به شدت بالا می‌برد.
\item معاینهٔ دوره‌ای چشم لازم است!
\item \en{ANA+}، با جنسیت \en{F}، درگیری کمتر مفاصل، درگیری بدون علامت مفاصل و سن پایین‌تر شروع همراه است.
\end{itemize}

\subsection*{کرایتریا تشخیص \en{JIA / JRA}}
\begin{itemize}
\item سن کمتر از ۱۶ سال.
\item آرتریت یک یا چند مفصل.
\item حداقل ۶ هفته.
\item تعیین \en{Onset Type} بیماری.
\item رد سایر علل آرتریت جوانان.
\end{itemize}
معیارها هر کدام از ساب‌تایپ‌های \en{JIA} در کتاب.

\subsection*{ویژگی‌های \en{sJIA}}
\begin{itemize}
\item خارج مفصلی: تب، راش ناپایدار، هپاتواسپلنومگالی، پلوریت.
\item \en{lab}: آنمی، \en{WBC\up\up}، \en{ESR\up\up}، \en{CRP\up\up}، \en{Ferritin\up}، \en{Plt\up\up}.
\item درمان: به \en{MTX} و \en{anti-TNF} کمتر پاسخ می‌دهند؛ معمولاً نیازمند اضافه کردن مهارکننده‌های \en{IL-1} و \en{IL-6}.
\end{itemize}

\clearpage
\pagehead{۳}
\subsection*{اولیگوآرتریت}
\begin{itemize}
\item الگوی درگیری مفاصل: زانو \en{++}، مچ پا، انگشتان \en{+}.
\item یووئیت در ۳۰\% موارد.
\item \en{ANA+} در ۶۰\% موارد؛ سایر آزمایش‌ها \en{NL} است، ممکن است \en{ESR/CRP} اندکی افزایش یافته.
\item درمان: \en{NSAID}، \en{CS} داخل مفصلی؛ ممکن است به \en{MTX} نیاز پیدا کند.
\end{itemize}

\subsection*{پلی‌آرتریت}
\textbf{\en{RF+}: تهاجمی و سیستمیک}
\begin{itemize}
\item در ۱۰\% بیماران ندول‌های روماتوئید، تب \en{low-grade}.
\item \en{ESR\up\up}، \en{CRP\up} یا \en{NL}، آنمی خفیف.
\item رمیشن طولانی‌مدت نامحتمل است؛ درمان سریع و تهاجمی.
\end{itemize}
\textbf{\en{RF$-$}:}
\begin{itemize}
\item قرینه یا غیرقرینه؛ مفاصل کوچک یا بزرگ؛ ستون مهرهٔ گردنی؛ \en{TMJ}.
\item نسبت \en{F/M: 9/1}.
\item یووئیت در ۱۰\%.
\item \en{ANA+} در ۴۰\% بیماران؛ \en{ESR\up} یا \en{\up\up}، \en{CRP NL} یا \en{\up}، آنمی خفیف.
\item \en{NSAID + MTX}؛ اگر پاسخ نداد: \en{anti-TNF (abatacept)}.
\end{itemize}

\subsection*{آرتریت پسوریاتیک}
\begin{itemize}
\item درگیری غیرقرینهٔ مفاصل کوچک تا متوسط.
\item \rng{2}{4} سال و \rng{9}{11} سال؛ \en{F/M = 2/1}.
\item یووئیت در ۱۰\%، پسوریازیس در ۵۰\%.
\item \en{ANA+} در ۵۰\%؛ \en{ESR\up}، \en{CRP\up} یا \en{NL}، آنمی خفیف.
\item \en{NSAID}، \en{CS} داخل مفصلی، \en{MTX}، \en{anti-TNF}.
\end{itemize}

\subsection*{\en{ERA} (آرتریت مرتبط با انتزیت)}
\begin{itemize}
\item سن شایع: \rng{9}{12} سال؛ \en{F/M = 1/7}.
\item بیشتر مفاصل اندام تحتانی و برخی مواقع اسکلت محوری.
\item یووئیت قدامی حاد و آنتریت ری‌اکتیو، \en{IBD}.
\item ۸۰\% بیماران \en{HLA-B27} مثبت هستند.
\item \en{NSAID}، \en{CS} داخل مفصلی، سولفاسالازین، گاهی \en{MTX}، \en{anti-TNF}.
\end{itemize}

\begin{itemize}
\item پلی‌آرتریت \en{RF} مثبت شبیه \en{RA} است: ندول‌های روماتوئید، سطوح اکستانسور آرنج، ستون مهره، تاندون آشیل.
\item درگیری مزمن مفصل \en{TMJ} منجر به \en{Micrognathia} می‌شود.
\end{itemize}

\clearpage
\pagehead{۴}
\begin{itemize}
\item درگیری ستون مهرهٔ سرویکال: اغلب به‌صورت کاهش \en{Extension} گردن است.
\item افزایش خطر ساب‌لاکسیشن آتلانتوآگزیال؛ علائم نورولوژیک.
\item درگیری هیپ: اغلب بی‌علامت یا خفیف.
\item \en{ROM} کاهش‌یافته یا دردناک.
\end{itemize}

\subsection*{تب در \en{sJIA}}
\begin{itemize}
\item اسپایک‌ها تب \en{$\geq 39$} درجه هر روزه یا هر دو روز که به‌صورت \en{$<37$} سقوط می‌کند و برای \rng{1}{2} هفته ادامه دارد.
\item اغلب عصرها رخ می‌دهد و با \sout{\en{Joint} و \en{Rash} همراه است}.
\item راش‌های پوستی همراه است.
\item روی تنه و پروگزیمال اندام‌ها: \en{Evanescent salmon-colored lesions}.
\item غیرخارش‌دار و ناپایدار است و زیر ۱ ساعت طول می‌کشد.
\item \en{Koebner phenomenon}: به دنبال تروما به پوست ضایعاتی ایجاد می‌شود.
\item تب، راش، هپاتواسپلنومگالی، لنفادنوپاتی در بیش از ۷۰\% بیماران دیده می‌شود.
\end{itemize}

\subsection*{سندرم فعال شدن ماکروفاژها (\en{MAS})}
\begin{itemize}
\item یک عارضهٔ نادر ولی کشنده در \en{sJIA} است.
\item در هر زمانی از بیماری ممکن است رخ بدهد.
\item نام دیگر آن \en{Secondary HLH} است.
\item شروع حاد تب‌های \en{high spiking}، لنفادنوپاتی، هپاتواسپلنومگالی، انسفالوپاتی.
\item \en{Plt\dn}، \en{WBC\dn}، \en{LFT\up}، \en{LDH\up}، \en{Ferritin\up}، \en{TG\up}.
\item \en{ESR} به علت هایپوفیبرینوژنمی پایین است.
\item برای افتراق \en{MAS} از \en{flare} بیماری زمینه‌ای.
\end{itemize}

\subsection*{کرایتریا تشخیص \en{sJIA-associated MAS}}
هایپرفریتینمی (\en{$>684$}) + حداقل ۲ عدد از موارد زیر:
\begin{itemize}
\item ترومبوسایتوپنی \en{$<181$}.
\item \en{AST $>48$}.
\item \en{TG $>156$}.
\item فیبرینوژن \en{$\leq 360$}.
\end{itemize}
نباید بیماری‌های زیر وجود داشته باشد:
\en{ITP}، هپاتیت عفونی، هایپر \en{TG} فامیلیال، لیشمانیوز احشایی.

\textbf{درمان اورژانسی:} متیل \en{high-dose}، سیکلوسپورین، آناکینرا.

\clearpage
\pagehead{۵}
\subsection*{معاینهٔ چشم‌پزشکی در بیمار \en{JIA}}
\begin{itemize}
\item در زمان تشخیص \en{JIA} همهٔ بیماران باید یک نوبت ویزیت شوند.
\item بیمار علامت‌دار چشمی حداکثر ظرف ۱ هفته؛ سایر موارد ظرف ۶ هفته.
\item بعد از شروع آرتریت هر ۲ ماه، بعد تا ۶ ماه، سپس هر ۴ ماه تا:
\end{itemize}

\textbf{الیگوآرتیکولر، پسوریاتیک، \en{ERH}، که زیر ۱۱ سال شروع شده‌اند:}
% ERH is the literal lettering on page 5; page 3 uses ERA.
\begin{center}
\begin{tabular}{|r|r|}
\hline
سن & مدت ویزیت چشم\\\hline
زیر ۳ سال & ۸ سال\\\hline
\rng{3}{4} سال & ۶ سال\\\hline
\rng{5}{8} سال & ۳ سال\\\hline
\rng{9}{10} سال & ۱ سال\\\hline
\end{tabular}
\end{center}

\textbf{پلی‌آرتیکولر، \en{JIA ANA+}، شروع زیر ۱۰ سال:}
\begin{center}
\begin{tabular}{|r|r|}
\hline
سن & مدت ویزیت چشم\\\hline
زیر ۶ سال & ۵ سال\\\hline
\rng{6}{9} سال & ۲ سال\\\hline
\end{tabular}
\end{center}

\begin{itemize}
\item پلی‌آرتیکولر، \en{ANA$-$}، شروع زیر ۷ سال: ۵ سال ویزیت چشم.
\item همهٔ گروه‌ها، سن بالای ۱۱ سال: ۱ سال ویزیت چشم.
\end{itemize}

\begin{itemize}
\item سطح \en{Hb} در \en{JIA} معمولاً \rng{7}{10} است. \en{ESR} معمولاً بالاست به‌جز در \en{MAS}.
\end{itemize}

\subsection*{درمان}
\begin{itemize}
\item هدف درمان: دستیابی به رمیشن، جلوگیری از آسیب مفصل، رشد نرمال کودک.
\item اولیگوآرتیکولار معمولاً به \en{NSAID} پاسخ می‌دهند و درد و التهاب کاهش می‌یابد.
\item تزریق \en{CS} داخل مفصلی در صورت عدم پاسخ به \en{NSAID} یا پاسخ نسبی (پس از \rng{4}{6} هفته)، محدودیت عملکردی و حرکتی، اختلاف طول دو اندام.
\item تریامسینولون.
\item درصد قابل توجهی از بیماران به این دو درمان پاسخ نمی‌دهند: \en{DMARD}.
\end{itemize}
\begin{enumerate}
\item \en{MTX}؛ \en{csDMARD}.
\item \en{anti-TNF}؛ \en{bDMARD}.
\end{enumerate}

\clearpage
\pagehead{۶}
\begin{itemize}
\item \rng{6}{12} هفته طول می‌کشد تا اثر \en{MTX} بروز کند.
\item عدم پاسخ به \en{MTX}: اضافه کردن \en{anti-TNF} (اتانرسپت؛ آدالیمومب).
\item از \en{Abatacept} (مهارکنندهٔ القای فعال شدن \en{T-cell})، \en{Tocilizumab} (آنتاگونیست گیرندهٔ \en{IL-6}) برای درمان فرم پلی‌آرتیکولار آن نیز استفاده می‌شود.
\item بیشترین خطر ایجاد یووئیت مزمن: اولیگوآرتریت خصوصاً \en{F}، \en{ANA+}، با شروع آرتریت زیر ۶ سال.
\end{itemize}

\clearpage
\section*{خلاصهٔ تکمیلی: \en{JIA -- 89}}
\pagehead{تکمیلی ۱ (تصویر ۳)}
% Stationery labels: Subject:, Year:, Month:, Date:; ruled-line numbers 1--9.
% Handwritten footer: 1. Blank date fields are retained below.
\noindent\en{Subject: JIA -- 89}\par
\en{Year: \ldots\quad Month: \ldots\quad Date: \ldots}

\subsection*{تعریف، تشخیص و دسته‌بندی}
\begin{itemize}
\item \en{JIA}: اختلالات مختلفی که منجر به آرتریت ماندگار در کودکان \textbf{زیر ۱۶ سال} می‌شوند.
\item تشخیص: آرتریت \en{$\geq 1$} مفصل، \textbf{مدت \en{$\geq 6$} هفته}، شروع زیر ۱۶ سال؛ رد سایر علل.
\item دسته‌بندی \en{JIA} بر اساس تعداد و ویژگی مفاصل درگیر \textbf{در ۶ ماه نخست بیماری}:
آرتریت سیستمیک (بیماری \en{Still})، آرتریت الیگوآرتیکولر، آرتریت پلی‌آرتیکولر، آرتریت مرتبط با انتزیت، آرتریت پسوریاتیک.
\end{itemize}

\subsection*{آرتریت سیستمیک (بیماری استیل)}
\begin{itemize}
\item در هر سنی زیر ۱۶ سال؛ \en{$M=F$}.
\item \textbf{حملات تب} (\en{$>38.5$}): \textbf{یک یا دو بار در روز، \en{$\geq 2$} بار در هفته}.
\item دورهٔ تب کوتاه است؛ در طول تب حال بیمار بد است.
\item علائم خارج مفصلی: \textbf{راش} ماکولوپاپولر اریتماتو؛ \textbf{لنفادنوپاتی}؛ \textbf{هپاتواسپلنومگالی}.
\item سروزیت؛ پریکاردیت؛ \textbf{میوکاردیت}.
\item لکوسیتوز؛ ترومبوسیتوز؛ آنمی؛ \en{ESR\up} و \en{CRP\up}.
\item ممکن است آرتریت با تأخیر چند ماهه بعد از سایر علائم بروز کند.
\end{itemize}

\subsection*{آرتریت الیگوآرتیکولر (\en{1--4} مفصل)}
\begin{itemize}
\item \textbf{شایع‌ترین تیپ \en{JIA}}؛ \textbf{زیر ۵ سال}؛ \en{$M<F$}.
\item اگر تعداد مفاصل درگیر ثابت باشد: \en{Persistent}؛ در غیر این صورت: \en{Extended}.
\item به‌طور شایع درگیری مفاصل \textbf{بزرگ}: \textbf{زانو (شایع‌ترین)}، مچ پا، مچ دست، آرنج.
\item \textbf{درگیری هیپ ناشایع است}.
\item در \rng{80}{90}\% موارد \textbf{\en{ANA} مثبت}؛ در همه \en{RF} منفی.
\item در \en{$1/3$} موارد \textbf{یووئیت قدامی}.
\item \en{Leg length discrepancy}.
\end{itemize}

\subsection*{آرتریت پلی‌آرتیکولر (\en{$\geq 5$} مفصل)}
\begin{itemize}
\item ۵۰\% \en{ANA} مثبت؛ ۱۰\% \textbf{یووئیت}.
\end{itemize}

\clearpage
\pagehead{تکمیلی ۲ (تصویر ۴)}
% Stationery labels: Subject:, Year:, Month:, Date:; ruled-line numbers 1--9.
% Handwritten footer: 2.
\noindent\en{Subject: JIA -- 89}\par
\en{Year: \ldots\quad Month: \ldots\quad Date: \ldots}

\subsection*{آرتریت پلی‌آرتیکولر: فرم \en{RF} منفی}
\begin{itemize}
\item دوم شایع؛ \textbf{\rng{2}{4} سال و \rng{9}{12} سال}؛ \en{$M<F$}.
\item درگیری سیمتریکال مفاصل بزرگ و کوچک دست و پا، \en{TMJ}، سرویکال اسپاین.
\end{itemize}

\subsection*{آرتریت پلی‌آرتیکولر: فرم \en{RF} مثبت}
\begin{itemize}
\item \textbf{اوایل نوجوانی}؛ \en{$M<F$}.
\item تهاجمی است و تا بزرگسالی طول می‌کشد و مانند \en{RA} می‌باشد.
\item به سرعت تخریب مفصلی می‌دهد؛ درگیری قرینهٔ مفاصل بزرگ و کوچک.
\item ندول‌های روماتوئید در آرنج و زانو.
\end{itemize}

\subsection*{آرتریت مرتبط با انتزیت}
\begin{itemize}
\item \textbf{نوجوانی}؛ \en{$F<M$}؛ \en{RF} منفی است.
\item مفاصل حامل وزن، خصوصاً هیپ؛ انتزیت آشیل؛ ساکروایلیاک.
\item می‌تواند در بزرگسالی تبدیل به \en{AS} شود؛ \textbf{مرتبط با \en{HLA-B27}}.
\item درگیری غیرقرینهٔ مفاصل اندام تحتانی؛ \textbf{یووئیت و آیریت}.
\end{itemize}

\subsection*{آرتریت پسوریاتیک}
\begin{itemize}
\item \rng{2}{4} سال و \rng{9}{11} سال؛ \en{$M<F$}.
\item آرتریت + پسوریازیس، یا آرتریت + حداقل ۲ مورد از: داکتیلیت، پیتینگ ناخن‌ها، سابقهٔ خانوادگی پسوریازیس.
\item درگیری مفاصل کوچک یا بزرگ (معمولاً زانو، مفاصل دست و پا).
\item ضایعات اریتماتو و \en{Scaly} در نواحی پست‌اوریکولار، پری‌آمبیلیکال و سطوح اکستانسور آرنج و زانو.
\end{itemize}

\subsection*{درمان}
\begin{itemize}
\item خط اول: \en{NSAID}، کورتیکواستروئید داخل مفصلی.
\item خط دوم: \en{DMARD}؛ \textbf{متوترکسات، سولفاسالازین، لفلونوماید}.
\item کورتیکواستروئید.
\item بیولوژیک: \textbf{\en{anti-TNF}} برای \textbf{پلی‌آرتیکولار}؛ اتانرسپت، اینفلیکسیمب، آدالیمومب.
\item \en{anti-IL-1} (آناکینرا)، \en{anti-IL-6} برای سیستمیک؛ توسیلیزومب (\en{anti-IL-6}).
\end{itemize}

\clearpage
\section*{تصاویر منبع}
\IfFileExists{images/pic01.jpg}{\def\imagebase{images/}}{%
\def\imagebase{pediatrics/ped-rheumatology/JIA/images/}}
\begin{center}
\textbf{تصویر 1}\par
\includegraphics[width=\linewidth,height=0.88\textheight,keepaspectratio]{\imagebase pic01.jpg}
\end{center}
\clearpage
\begin{center}
\textbf{تصویر 2}\par
\includegraphics[width=\linewidth,height=0.88\textheight,keepaspectratio]{\imagebase pic02.jpg}
\end{center}
\clearpage
\begin{center}
\textbf{تصویر ۳: خلاصهٔ تکمیلی، صفحهٔ ۱}\par
\includegraphics[width=\linewidth,height=0.88\textheight,keepaspectratio]{\imagebase pic03.jpg}
\end{center}
\clearpage
\begin{center}
\textbf{تصویر ۴: خلاصهٔ تکمیلی، صفحهٔ ۲}\par
\includegraphics[width=\linewidth,height=0.88\textheight,keepaspectratio]{\imagebase pic04.jpg}
\end{center}
\clearpage
\begin{center}
\textbf{تصویر 5}\par
\includegraphics[width=\linewidth,height=0.88\textheight,keepaspectratio]{\imagebase pic05.jpg}
\end{center}
\clearpage
\begin{center}
\textbf{تصویر 6}\par
\includegraphics[width=\linewidth,height=0.88\textheight,keepaspectratio]{\imagebase pic06.jpg}
\end{center}
\end{document}
